\subsection{Single-bond peeling and the window witness}
\label{subsec:peps_torus_single_bond_peeling_window_witness}

\begin{definition}[Reference edge of the horizontal staircase]
    \label{def:peps_horizontal_staircase_reference_edge}
    \lean{TNLean.PEPS.horizontalStaircaseReferenceEdge}
    \leanok
    \uses{def:peps_horizontal_staircase_patch}
    In the staircase coordinates $s=(a,b)$, the reference edge is the
    horizontal edge joining $(a+L-1,b+K-1)$ to $(a+L,b+K-1)$.
    It is the edge between the left and right end windows.
\end{definition}

\begin{theorem}[The reference edge is the unique crossing of the two end windows]
    \label{thm:peps_horizontal_staircase_reference_edge_crossing}
    \lean{TNLean.PEPS.isCrossingEdge_horizontalStaircaseEndWindows}
    \lean{TNLean.PEPS.isRegionBoundaryEdge_horizontalStaircaseLeftWindow_referenceEdge}
    \lean{TNLean.PEPS.isRegionBoundaryEdge_horizontalStaircaseRightWindow_referenceEdge}
    \lean{TNLean.PEPS.not_isRegionBoundaryEdge_horizontalStaircaseEndPair_referenceEdge}
    \leanok
    \uses{def:peps_horizontal_staircase_reference_edge, def:peps_horizontal_staircase_patch}
    Assume
    $1<W$, $1<H$, $0<L$, $0<K$, $1\le a$, $a+2L\le W$, and
    $b+2K-1\le H$.
    Then an edge crosses from the left end window to the right end window if
    and only if it is the reference edge $e$:
    \begin{align}
        f\in\partial W_{\mathrm{left}}\cap\partial W_{\mathrm{right}}
        \quad\Longleftrightarrow\quad
        f=e.
        \notag
    \end{align}
    Thus $e$ is a boundary edge of each end window, but both endpoints of $e$
    lie in the end pair $S$, so $e\notin\partial S$.
\end{theorem}
\begin{proof}
    \leanok
    The endpoint calculation identifies $e$ with the unique horizontal edge
    joining the two displayed rectangles. Since one endpoint lies in
    $W_{\mathrm{left}}$ and the other lies in $W_{\mathrm{right}}$, the same
    edge is internal to their union $S$.
\end{proof}

\begin{lemma}[Boundary edges of the end pair come from the end windows]
    \label{lem:peps_horizontal_staircase_end_pair_boundary_from_windows}
    \lean{TNLean.PEPS.isRegionBoundaryEdge_endPair_of_leftWindow}
    \lean{TNLean.PEPS.isRegionBoundaryEdge_endPair_of_rightWindow}
    \lean{TNLean.PEPS.isRegionBoundaryEdge_window_of_endPair}
    \lean{TNLean.PEPS.ne_referenceEdge_of_isRegionBoundaryEdge_endPair}
    \leanok
    \uses{thm:peps_horizontal_staircase_reference_edge_crossing}
    Under the bounds
    $1<W$, $1<H$, $0<L$, $0<K$, $1\le a$, $a+2L\le W$, and
    $b+2K-1\le H$,
    a boundary edge of the end pair is a boundary edge of one of the two end
    windows, and it is not the reference edge. Conversely, any boundary edge
    of one end window other than $e$ is a boundary edge of $S$.
\end{lemma}
\begin{proof}
    \leanok
    A boundary edge of $S$ has one endpoint in $S$; that endpoint lies in one
    of the two end windows. The other endpoint is outside $S$, hence outside
    that window. Conversely, a boundary edge of one window which is not $e$
    cannot enter the other window, because $e$ is the unique crossing edge.
\end{proof}

\begin{theorem}[Boundary of the staircase end pair]
    \label{thm:peps_horizontal_staircase_end_pair_boundary}
    \lean{TNLean.PEPS.eq_referenceEdge_of_isRegionBoundaryEdge_both_endWindows}
    \lean{TNLean.PEPS.isRegionBoundaryEdge_endPair_iff}
    \lean{TNLean.PEPS.isRegionBoundaryEdge_endPair_exactlyOne_window}
    \leanok
    \uses{lem:peps_horizontal_staircase_end_pair_boundary_from_windows}
    Under the bounds
    $1<W$, $1<H$, $0<L$, $0<K$, $1\le a$, $a+2L\le W$, and
    $b+2K-1\le H$,
    an edge is a boundary edge of the staircase end pair $S$ if and only if it
    is a boundary edge of one of the two end windows and is not the reference
    edge:
    \begin{align}
        f\in\partial S
        \quad\Longleftrightarrow\quad
        \bigl(f\in\partial W_{\mathrm{left}}
        \quad\text{or}\quad
        f\in\partial W_{\mathrm{right}}\bigr)
        \ \text{and}\ f\ne e.
        \notag
    \end{align}
    Moreover each edge in $\partial S$ lies on exactly one of the two window
    boundaries.
\end{theorem}
\begin{proof}
    \leanok
    Lemma~\ref{lem:peps_horizontal_staircase_end_pair_boundary_from_windows}
    gives both implications. If an edge lay on both end-window boundaries,
    the unique-crossing theorem would identify it with $e$, which is not a
    boundary edge of $S$.
\end{proof}

\begin{theorem}[Boundary configurations of the staircase end windows]
    \label{thm:peps_horizontal_staircase_end_window_boundary_config}
    \lean{TNLean.PEPS.horizontalStaircaseLeftWindowBoundaryConfigEquiv}
    \lean{TNLean.PEPS.horizontalStaircaseRightWindowBoundaryConfigEquiv}
    \lean{TNLean.PEPS.horizontalStaircaseEndPairBoundaryConfigEquiv}
    \leanok
    \uses{thm:peps_horizontal_staircase_end_pair_boundary, thm:peps_horizontal_staircase_reference_edge_crossing}
    Under the same geometric bounds, write
    $\partial_{\mathrm{left}}S$ and $\partial_{\mathrm{right}}S$
    for the boundary edges of $S$ which lie on the corresponding end-window
    boundary.  Then the boundary-configuration spaces have canonical
    factorizations
    \begin{align}
        \operatorname{Conf}(\partial W_{\mathrm{left}})
          &\cong \operatorname{Conf}(e)
            \times\operatorname{Conf}(\partial_{\mathrm{left}}S),
        \notag \\
        \operatorname{Conf}(\partial W_{\mathrm{right}})
          &\cong \operatorname{Conf}(e)
            \times\operatorname{Conf}(\partial_{\mathrm{right}}S),
        \notag \\
        \operatorname{Conf}(\partial S)
          &\cong
            \operatorname{Conf}(\partial_{\mathrm{left}}S)
              \times\operatorname{Conf}(\partial_{\mathrm{right}}S).
        \notag
    \end{align}
    These equivalences preserve every underlying lattice-edge label.  Thus
    the reference edge is the only contracted virtual index shared by the
    two highlighted end windows; every other virtual index remains on the
    corresponding open boundary of $S$.
\end{theorem}
\begin{proof}
    \leanok
    Split each end-window configuration at the reference edge.  A remaining
    boundary edge of either window is a boundary edge of $S$, and conversely
    every boundary edge of $S$ belongs to exactly one end window.  Reindexing
    the dependent products along these edge identifications gives the three
    displayed equivalences.
\end{proof}

\begin{lemma}[Complement split for an end window]
    \label{lem:peps_horizontal_staircase_end_window_complement_split}
    \lean{TNLean.PEPS.horizontalStaircaseEndPair_sdiff_leftWindow}
    \lean{TNLean.PEPS.horizontalStaircaseEndPair_sdiff_rightWindow}
    \lean{TNLean.PEPS.compl_leftWindow_eq_rightWindow_union_complEndPair}
    \lean{TNLean.PEPS.compl_rightWindow_eq_leftWindow_union_complEndPair}
    \leanok
    \uses{def:peps_horizontal_staircase_patch, lem:peps_horizontal_staircase_end_pair_disjoint}
    If $2L\le W$, the complement of either end window splits as the opposite
    end window and the complement of the end pair:
    \begin{align}
        V\setminus W_{\mathrm{left}}
         & =W_{\mathrm{right}}\cup (V\setminus S),
        \notag                                     \\
        V\setminus W_{\mathrm{right}}
         & =W_{\mathrm{left}}\cup (V\setminus S).
        \notag
    \end{align}
\end{lemma}
\begin{proof}
    \leanok
    Since the two end windows are disjoint and
    $S=W_{\mathrm{left}}\cup W_{\mathrm{right}}$, removing one window from
    $S$ leaves the other. Splitting $V\setminus W$ through the inclusion
    $W\subseteq S$ gives the displayed decompositions.
\end{proof}

\begin{theorem}[The reference edge bounds the two end windows]
    \label{thm:peps_staircase_end_reference_edge_boundary}
    \lean{TNLean.PEPS.isRegionBoundaryEdge_staircaseWindow_zero_referenceEdge}
    \lean{TNLean.PEPS.isRegionBoundaryEdge_staircaseWindow_last_referenceEdge}
    \leanok
    Let the torus have width $W$ and height $H$, with $1<H$.
    Let $e$ be the horizontal reference edge of a staircase window family with
    dimensions $L,K$ and lower-left coordinate $(a,b)$. Suppose
    $0<L$, $0<K$, $1\le a$, $a+2L\le W$, and $b+2K-1\le H$.
    Then $e$ is a boundary edge of the first staircase window $W_0$ and of the
    last staircase window $W_{L+K-1}$:
    \begin{align}
        e & \in\partial W_0,
        \notag                     \\
        e & \in\partial W_{L+K-1}.
        \notag
    \end{align}
\end{theorem}
\begin{proof}
    \leanok
    The first staircase window is the right end window, and the last
    staircase window is the left end window:
    \begin{align}
        W_0       & =W_{\mathrm{right}},
        \notag                           \\
        W_{L+K-1} & =W_{\mathrm{left}}.
        \notag
    \end{align}
    The reference edge is a boundary edge of these two end windows by their
    defining geometry.
\end{proof}

\begin{theorem}[Open-boundary equality on the staircase end pair]
    \label{thm:peps_staircase_pair_insert_eq_open}
    \lean{TNLean.PEPS.staircasePair_insert_eq_open}
    \leanok
    \uses{thm:peps_staircase_pair_cancel, thm:peps_corner_extension_linear_transitive}
    Let $S=W_0\cup W_{L+K-1}$ be the staircase end pair. Suppose that the
    one-orientation $L\times K$ window injectivity hypotheses hold for a
    tensor $B$, unions of injective regions are injective, all bond dimensions
    of $B$ are positive, and
    $2\le L$, $2\le K$, $2L+1\le W$, and $2K+1\le H$.
    If $C_j$ are inserts on the staircase windows whose deformed states all
    equal one common state, then the two end-window inserts have equal corner
    extensions to $S$:
    \begin{align}
        \Ext_{W_0\subseteq S}(C_0)
        =
        \Ext_{W_{L+K-1}\subseteq S}(C_{L+K-1}).
        \notag
    \end{align}
\end{theorem}
\begin{proof}
    \leanok
    The consecutive-window equalities are composed across the staircase patch.
    Both sides are then extended through the completed corner, and
    transitivity of corner extension rewrites the two extension chains. The
    completed corner is an injective $L\times K$ window, so
    Theorem~\ref{thm:peps_staircase_pair_cancel} cancels it and leaves the
    equality on the end pair $S$.
\end{proof}

\begin{theorem}[Lemma 5 comparison of the two end blocks]
    \label{thm:peps_existsUnique_virtualOperation_of_endPair}
    \lean{TNLean.PEPS.existsUnique_virtualOperation_of_endPair}
    \leanok
    Let $K$ be the basis of the virtual space on the bond between two
    injective blocks, and let
    $\{A_{1,\mu}\}_{\mu\in K}$ and
    $\{A_{2,\mu}\}_{\mu\in K}$ be linearly independent families.  Suppose
    that the two end physical operations $O_1$ and $O_3$ produce families
    $C_1$ and $C_2$ satisfying the open contraction identity
    \begin{align}
        \sum_{\mu\in K} C_{1,\mu}\otimes A_{2,\mu}
        =
        \sum_{\nu\in K} A_{1,\nu}\otimes C_{2,\nu}.
        \notag
    \end{align}
    Then there is a unique matrix $X\in\MN{|K|}$ such that
    \begin{align}
        C_{1,\nu}
        &=\sum_{\mu\in K} A_{1,\mu}X_{\mu,\nu},
        &
        C_{2,\mu}
        &=\sum_{\nu\in K} X_{\mu,\nu}A_{2,\nu}.
        \notag
    \end{align}
    Thus the two end operations are obtained by inserting the same virtual
    operation $X$ on their common bond.  No injectivity is required of the
    modified families $C_1$ and $C_2$.  This is the ``compare the two ends''
    step in Lemma~5 of \cite{Molnar2018NormalPEPS}.
\end{theorem}
\begin{proof}
    \leanok
    Apply a dual functional to the linearly independent family
    $\{A_{2,\mu}\}$ in the open contraction identity.  This expresses every
    $C_{1,\nu}$ as a linear combination of the $A_{1,\mu}$ with a coefficient
    matrix $X$.  Substitution into the same identity and linear independence
    of $\{A_{1,\mu}\}$ gives the second displayed relation with the same
    matrix.  Linear independence of the first genuine family also makes the
    coefficient matrix unique.
\end{proof}

\begin{theorem}[Single-bond peeling of the staircase end equality]
    \label{thm:peps_regionInsertedCoeff_endWindows_eq_of_staircase}
    \lean{TNLean.PEPS.regionInsertedCoeff_endWindows_eq_of_staircase}
    \leanok
    \uses{def:peps_bond_inserted_region_insert, thm:peps_regionInsertedCoeff_eq_extendInsert_bondInserted, thm:peps_staircase_end_reference_edge_boundary, thm:peps_staircase_pair_insert_eq_open}
    Let the torus have width $W$ and height $H$.
    Suppose that the one-orientation $L\times K$ window injectivity
    hypotheses hold for a tensor $B$, unions of injective regions are
    injective, and all bond dimensions of $B$ are positive. Assume
    \begin{align}
        2\le L,\qquad 2\le K,\qquad 1\le a,\qquad
        a+2L\le W,\qquad b+2K-1\le H,
        \notag \\
        2L+1\le W,\qquad 2K+1\le H.
        \notag
    \end{align}
    Let $C_j$ be inserts on the staircase windows $W_j$ whose deformed states
    all equal one common state. Suppose that
    $C_0=I_{B,W_0,e,M}$ and
    $C_{L+K-1}=I_{B,W_{L+K-1},e,M'}$.
    Then for every global physical configuration $\sigma$,
    \begin{align}
        C_B(W_0,e,M;\sigma|_{W_0},\sigma|_{V\setminus W_0})
         & =
        C_B\bigl(W_{L+K-1},e,M';
        \sigma|_{W_{L+K-1}},\sigma|_{V\setminus W_{L+K-1}}\bigr).
        \notag
    \end{align}
\end{theorem}
\begin{proof}
    \leanok
    Theorem~\ref{thm:peps_regionInsertedCoeff_eq_extendInsert_bondInserted}
    rewrites the two coefficients as the deformed states on the staircase end
    pair of the two corner-extended inserts:
    \begin{align}
        C_B(W_0,e,M;\sigma|_{W_0},\sigma|_{V\setminus W_0})
         & =
        \Psi_{B,S,\Ext_{W_0\subseteq S}(C_0)}(\sigma),
        \notag \\
        C_B(W_{L+K-1},e,M';
        \sigma|_{W_{L+K-1}},\sigma|_{V\setminus W_{L+K-1}})
         & =
        \Psi_{B,S,
            \Ext_{W_{L+K-1}\subseteq S}(C_{L+K-1})}(\sigma).
        \notag
    \end{align}
    The open-boundary equality on the staircase end pair gives
    $\Ext_{W_0\subseteq S}(C_0)
        =\Ext_{W_{L+K-1}\subseteq S}(C_{L+K-1})$, and hence the two coefficients
    are equal.
\end{proof}

\paragraph{Window-independence of the bond insertion.}

\begin{lemma}[Restriction and reassembly of a physical configuration]
    \label{lem:peps_assembleRegionSigma_restrict}
    \lean{TNLean.PEPS.assembleRegionσ_restrict}
    \leanok
    Let $\omega$ be a physical configuration on the full vertex set.  If
    $\omega_R$ and $\omega_{V\setminus R}$ are its restrictions to $R$ and
    to the complement, then
    \[
        \operatorname{Assemble}_R(\omega_R,\omega_{V\setminus R})=\omega.
    \]
\end{lemma}
\begin{proof}
    \leanok
    At a vertex in $R$ the assembled configuration reads $\omega_R$, and at a
    vertex outside $R$ it reads $\omega_{V\setminus R}$.  These restrictions
    are both inherited from $\omega$, so the two configurations agree at every
    vertex.
\end{proof}

\begin{theorem}[Bond-inserted state as an edge insertion]
    \label{thm:peps_deformedRegionStateAssembled_bondInserted_eq_smul_edgeInsertedCoeff}
    \lean{TNLean.PEPS.deformedRegionStateAssembled_bondInserted_eq_smul_edgeInsertedCoeff}
    \leanok
    \uses{thm:peps_deformedRegionState_bondInsertedRegionInsert,
        thm:peps_regionInsertedCoeff_eq_smul_edgeInsertedCoeff,
        lem:peps_assembleRegionSigma_restrict}
    Let $R$ be a region, let $f$ be a boundary edge of $R$, let
    $M\in\MN{D_B(f)}$, and let $\omega$ be a physical configuration on the
    full vertex set.  Then the assembled deformed state of the bond-inserted
    region insert is the non-boundary bond product times the corresponding
    edge-inserted coefficient:
    \[
        \Psi^{\mathrm{as}}_{B,R,I_{B,R,f,M}}(\omega)
        =
        p_B(R)\,
        E_B\bigl(f,\omega,\mathcal O_{R,f}(M)\bigr).
    \]
    Here $p_B(R)=\prod_{g\notin\partial R}D_B(g)$, and $\mathcal O_{R,f}$ is
    the boundary-edge orientation.
\end{theorem}
\begin{proof}
    \leanok
    Theorem~\ref{thm:peps_deformedRegionState_bondInsertedRegionInsert}
    rewrites the assembled deformed state as
    \[
        C_B(R,f,M;\omega_R,\omega_{V\setminus R}).
    \]
    Theorem~\ref{thm:peps_regionInsertedCoeff_eq_smul_edgeInsertedCoeff}
    turns this coefficient into
    \[
        p_B(R)\,
        E_B\bigl(f,\operatorname{Assemble}_R
            (\omega_R,\omega_{V\setminus R}),\mathcal O_{R,f}(M)\bigr).
    \]
    Lemma~\ref{lem:peps_assembleRegionSigma_restrict} identifies the assembled
    configuration with $\omega$.
\end{proof}

\begin{theorem}[Window-independence of a bond insertion]
    \label{thm:peps_bondInserted_windowIndependent}
    \lean{TNLean.PEPS.bondInserted_windowIndependent}
    \leanok
    \uses{thm:peps_deformedRegionStateAssembled_bondInserted_eq_smul_edgeInsertedCoeff}
    Let $R_1$ and $R_2$ be two regions for which the same edge $e$ is a
    boundary edge.  Let $M_1,M_2\in\MN{D_B(e)}$, and suppose that the two
    oriented insertions agree,
    \[
        \mathcal O_{R_1,e}(M_1)=\mathcal O_{R_2,e}(M_2).
    \]
    Then for every physical configuration $\omega$ on the full vertex set,
    \[
    \begin{aligned}
        p_B(R_2)\,
        \Psi^{\mathrm{as}}_{B,R_1,I_{B,R_1,e,M_1}}(\omega)
        &=
        p_B(R_1)\,
        \Psi^{\mathrm{as}}_{B,R_2,I_{B,R_2,e,M_2}}(\omega).
    \end{aligned}
    \]
\end{theorem}
\begin{proof}
    \leanok
    Applying
    Theorem~\ref{thm:peps_deformedRegionStateAssembled_bondInserted_eq_smul_edgeInsertedCoeff}
    to each of the two regions gives, for $i=1,2$,
    \[
        \Psi^{\mathrm{as}}_{B,R_i,I_{B,R_i,e,M_i}}(\omega)
        =
        p_B(R_i)\,E_B(e,\omega,\mathcal O_{R_i,e}(M_i)).
    \]
    The oriented insertions are equal by hypothesis, so both sides of the
    claimed identity are
    \[
        p_B(R_1)p_B(R_2)\,E_B(e,\omega,\mathcal O_{R_1,e}(M_1)).
    \]
\end{proof}

\paragraph{One virtual operation on every staircase window.}

The construction below formalizes the forward virtual/physical-operation
correspondence in the proof sketch of the two-dimensional corollary
\cite[corollary after Lemma~5, lines~2297--2444]{Molnar2018NormalPEPS}.

\begin{definition}[Partial state across a region cut]%
    \label{def:peps_region_partial_state}
    \lean{TNLean.PEPS.regionPartialState}
    \leanok
    For a region $R$ and a physical configuration $\tau$ on its complement,
    define the partial state $\psi^\tau_{A,R}$ on the physical space of $R$
    by
    \begin{align}
        \psi^\tau_{A,R}(\sigma)
         & =c_A(\omega_{R,\sigma,\tau}).
        \notag
    \end{align}
\end{definition}

\begin{theorem}[Partial states of equal closed states]%
    \label{thm:peps_region_partial_state_same_state}
    \lean{TNLean.PEPS.regionPartialState_sameState}
    \leanok
    \uses{def:peps_region_partial_state}
    If $A$ and $B$ generate the same closed state, then, for every region
    $R$ and every complement configuration $\tau$,
    $\psi^\tau_{A,R}=\psi^\tau_{B,R}$.
\end{theorem}
\begin{proof}
    \leanok
    Evaluate both functions at a physical configuration $\sigma$ on $R$ and
    apply the equality of the two closed-state coefficients at
    $\omega_{R,\sigma,\tau}$.
\end{proof}

\begin{theorem}[Blocked-tensor expansion of a partial state]%
    \label{thm:peps_region_partial_state_blocked_expansion}
    \lean{TNLean.PEPS.regionInteriorBondProd_smul_regionPartialState}
    \leanok
    \uses{def:peps_region_partial_state,
        def:peps_regionComplementBoundaryConfig}
    Let $P_A(R)$ be the product of the dimensions of the bonds internal to
    $R$. For every complement configuration $\tau$,
    \begin{align}
        P_A(R)\psi^\tau_{A,R}
         & =\sum_{\mu\in\partial R}
        T_{A,V\setminus R}(\bar\mu)(\tau)T_{A,R}(\mu).
        \label{eq:peps_partial_state_blocked_expansion}
    \end{align}
\end{theorem}
\begin{proof}
    \leanok
    Split the closed contraction across $R$. If $\mu$ and $\nu$ are boundary
    labels on $R$, the complement-boundary identification satisfies
    \begin{align}
        \bar\mu=\bar\nu
         & \quad\Longleftrightarrow\quad \mu=\nu.
        \notag
    \end{align}
    Hence the double boundary sum collapses to its diagonal, which gives the
    sum in~(\ref{eq:peps_partial_state_blocked_expansion}); each bond internal
    to $R$ contributes its dimension to the factor $P_A(R)$.
\end{proof}

\begin{definition}[Physical operator of a region insert]%
    \label{def:peps_physical_operator_of_region_insert}
    \lean{TNLean.PEPS.physicalOpOfRegionInsert}
    \leanok
    \uses{def:peps_region_out_endpoint_block_inverse}
    Let $R$ be an injective region and let
    $C=\{C(\mu)\}_{\mu\in\partial R}$ be any family in its physical
    space, indexed by the virtual boundary configurations.  If $L_{A,R}$
    is a left inverse of the blocked tensor map, define
    \begin{align}
        O_C
         & =\left(c\longmapsto\sum_{\mu}c_\mu C(\mu)\right)L_{A,R}.
        \notag
    \end{align}
\end{definition}

\begin{theorem}[Open-boundary realization of a region insert]%
    \label{thm:peps_physical_operator_of_region_insert_realizes}
    \lean{TNLean.PEPS.physicalOpOfRegionInsert_realizes}
    \leanok
    \uses{def:peps_physical_operator_of_region_insert}
    For every virtual boundary configuration $\mu$,
    \begin{align}
        O_C\bigl(T_{A,R}(\mu)\bigr) & =C(\mu).
        \notag
    \end{align}
    Thus the insert is obtained from the genuine region block by one
    physical operator before the region is contracted with its complement.
\end{theorem}
\begin{proof}
    \leanok
    The left inverse sends $T_{A,R}(\mu)$ to the standard basis vector
    indexed by $\mu$, and the displayed linear combination sends that basis
    vector to $C(\mu)$.
\end{proof}

\begin{definition}[Region insert of a physical operator]%
    \label{def:peps_region_insert_of_physical_operator}
    \lean{TNLean.PEPS.regionInsertOfPhysicalOp}
    \leanok
    Let $O$ be a physical operator on the physical space of a region $R$.
    Its action on the genuine open-boundary blocks defines the insert
    \begin{align}
        C^O_{A,R}(\mu) & =O\bigl(T_{A,R}(\mu)\bigr).
        \notag
    \end{align}
\end{definition}

\begin{theorem}[Closed state of a physical operator]%
    \label{thm:peps_closed_state_of_physical_operator}
    \lean{TNLean.PEPS.deformedRegionState_regionInsertOfPhysicalOp}
    \leanok
    \uses{def:peps_region_insert_of_physical_operator,
        thm:peps_region_partial_state_blocked_expansion}
    For a complement physical configuration $\tau$, let
    $\psi^\tau_{A,R}$ be the partial state across the cut. Then
    \begin{align}
        \Psi_{A,R,C^O_{A,R}}(\sigma,\tau)
         & =\bigl[O\bigl(P_A(R)\psi^\tau_{A,R}\bigr)\bigr](\sigma).
        \label{eq:peps_physical_op_partial_state}
    \end{align}
\end{theorem}
\begin{proof}
    \leanok
    Apply $O$ to~(\ref{eq:peps_partial_state_blocked_expansion}) and evaluate
    at $\sigma$. This proves
    (\ref{eq:peps_physical_op_partial_state}).
\end{proof}

\begin{theorem}[Physical operator transferred between equal closed states]%
    \label{thm:peps_physical_operator_same_state_transfer}
    \lean{TNLean.PEPS.deformedRegionState_regionInsertOfPhysicalOp_sameState}
    \leanok
    \uses{thm:peps_closed_state_of_physical_operator,
        thm:peps_region_partial_state_same_state}
    If $A$ and $B$ generate the same closed state, then
    \begin{align}
        P_B(R)\Psi_{A,R,C^O_{A,R}}(\sigma,\tau)
         & =P_A(R)\Psi_{B,R,C^O_{B,R}}(\sigma,\tau).
        \label{eq:peps_physical_op_same_state}
    \end{align}
    The two virtual boundary spaces need not be identified.
\end{theorem}
\begin{proof}
    \leanok
    Equality of the closed states gives
    $\psi^\tau_{A,R}=\psi^\tau_{B,R}$. Apply
    (\ref{eq:peps_physical_op_partial_state}) to both tensors and
    cross-multiply the two interior-bond products.
\end{proof}

\begin{definition}[Interior-bond region insert]%
    \label{def:peps_interior_bond_inserted_region_insert}
    \lean{TNLean.PEPS.interiorBondInsertedRegionInsert}
    \leanok
    \uses{def:peps_bond_inserted_region_insert, def:peps_corner_extension}
    Let $e=\{u,w\}$ be an ordered edge whose two endpoints lie in a region
    $R$, with $u$ its ordered left endpoint, and let
    $X\in\MN{D_A(e)}$.  Write $U=\{u\}$ and let $I_{A,U,e,X}$ be the
    boundary-bond insert on $U$.  The interior-bond insert on $R$ is
    \begin{align}
        I^{\mathrm{int}}_{A,R,e,X}
         & =
        \frac{P_A(R)}{P_A(U)}
        \Ext_{U\subseteq R}(I_{A,U,e,X}).
        \notag
    \end{align}
    Thus the genuine tensors at the vertices of $R\setminus U$ are adjoined
    to a boundary insertion of $X$ on $e$, with the normalization changed
    from the non-boundary bond product of $U$ to that of $R$.
\end{definition}

\begin{theorem}[Closed state of an interior-bond insert]%
    \label{thm:peps_interior_bond_inserted_state}
    \lean{TNLean.PEPS.deformedRegionStateAssembled_interiorBondInserted_eq_smul_edgeInsertedCoeff}
    \leanok
    \uses{def:peps_interior_bond_inserted_region_insert, thm:peps_corner_extension_state_preserves, thm:peps_deformedRegionStateAssembled_bondInserted_eq_smul_edgeInsertedCoeff}
    Suppose all bond dimensions are positive.  For every global physical
    configuration $\omega$,
    \begin{align}
        \Psi^{\mathrm{as}}_{A,R,I^{\mathrm{int}}_{A,R,e,X}}(\omega)
         & =P_A(R)E_A(e,\omega,X).
        \notag
    \end{align}
\end{theorem}
\begin{proof}
    \leanok
    Corner extension preserves the closed state.  On the singleton $U$, the
    boundary orientation is the identity because $U$ contains the ordered
    left endpoint of $e$.  Hence its boundary insert has assembled state
    $P_A(U)E_A(e,\omega,X)$.  The factor $P_A(R)/P_A(U)$ cancels $P_A(U)$
    and gives the displayed identity.
\end{proof}

\begin{definition}[Staircase family for one virtual operation]%
    \label{def:peps_staircase_virtual_operation_insert}
    \lean{TNLean.PEPS.staircaseVirtualOperationInsert}
    \leanok
    \uses{def:peps_bond_inserted_region_insert, def:peps_interior_bond_inserted_region_insert, lem:peps_reference_edge_boundary_windows, lem:peps_reference_edge_interior_windows}
    Let $e$ be the horizontal reference edge and let
    $X\in\MN{D_B(e)}$.  Define an insert $C_j(X)$ on each staircase window
    $W_j$, for $0\le j<L+K$, by
    \begin{align}
        C_j(X)
         & =
        \begin{cases}
            I_{B,W_0,e,X^{\mathsf T}},    & j=0,        \\
            I^{\mathrm{int}}_{B,W_j,e,X}, & 1\le j<L,   \\
            I_{B,W_j,e,X},                & L\le j<L+K.
        \end{cases}
        \notag
    \end{align}
    The first window contains only the ordered right endpoint of $e$, the
    middle windows contain both endpoints, and the remaining windows contain
    only the ordered left endpoint.  The transpose in the first case makes
    all three boundary orientations represent the same matrix $X$.
\end{definition}

\begin{theorem}[Physical operators on the staircase windows]%
    \label{thm:peps_staircase_virtual_operation_physical_realization}
    \lean{TNLean.PEPS.staircaseVirtualOperationPhysicalOp}
    \lean{TNLean.PEPS.staircaseVirtualOperationPhysicalOp_realizes}
    \leanok
    \uses{def:peps_staircase_virtual_operation_insert, def:peps_physical_operator_of_region_insert, thm:peps_physical_operator_of_region_insert_realizes}
    Suppose every $L\times K$ window is injective.  For every
    $0\leq j<L+K$, there is a physical operator $O_j(X)$ on the physical
    space of $W_j$ such that
    \begin{align}
        O_j(X)\bigl(T_{B,W_j}(\mu)\bigr) & =C_j(X)(\mu)
        \qquad (\mu\in\partial W_j).
        \notag
    \end{align}
    Hence the virtual operation is realized locally on each window, before
    any contraction with the complementary tensor network.
\end{theorem}
\begin{proof}
    \leanok
    Each staircase window is an $L\times K$ translate and is therefore
    injective.  Apply the open-boundary realization theorem to $C_j(X)$.
\end{proof}

\begin{lemma}[The two end inserts of the staircase family]%
    \label{lem:peps_staircase_virtual_operation_end_inserts}
    \lean{TNLean.PEPS.staircaseVirtualOperationInsert_zero}
    \lean{TNLean.PEPS.staircaseVirtualOperationInsert_last}
    \leanok
    \uses{def:peps_staircase_virtual_operation_insert}
    The end members of the family are
    \begin{align}
        C_0(X)       & =I_{B,W_0,e,X^{\mathsf T}},
                     &
        C_{L+K-1}(X) & =I_{B,W_{L+K-1},e,X}.
        \notag
    \end{align}
\end{lemma}
\begin{proof}
    \leanok
    The first identity is the $j=0$ case of the definition.  Since $K>0$,
    the last index satisfies $L\le L+K-1$, which gives the second identity.
\end{proof}

\begin{theorem}[Common state of the staircase virtual-operation family]%
    \label{thm:peps_staircase_virtual_operation_common_state}
    \lean{TNLean.PEPS.deformedRegionStateAssembled_staircaseVirtualOperationInsert}
    \leanok
    \uses{def:peps_staircase_virtual_operation_insert, lem:peps_staircase_virtual_operation_end_inserts, lem:peps_reference_edge_left_endpoint_staircase_window, thm:peps_interior_bond_inserted_state, thm:peps_deformedRegionStateAssembled_bondInserted_eq_smul_edgeInsertedCoeff, thm:peps_regionInteriorBondProd_staircaseWindow_eq}
    Let $B$ be translation invariant and suppose all its bond dimensions are
    positive.  Then for $0\le j<L+K$ and every global physical configuration
    $\omega$,
    \begin{align}
        \Psi^{\mathrm{as}}_{B,W_j,C_j(X)}(\omega)
         & =P_B(W_0)E_B(e,\omega,X).
        \notag
    \end{align}
    In particular, all $L+K$ inserts have one common deformed state.
\end{theorem}
\begin{proof}
    \leanok
    On $W_0$, the boundary orientation sends $X^{\mathsf T}$ to $X$.
    On $W_1,\ldots,W_{L-1}$, apply the interior-bond identity.  On the
    remaining windows the ordered left endpoint lies in the window, so the
    boundary orientation fixes $X$.  Each case gives
    $P_B(W_j)E_B(e,\omega,X)$.  Translation invariance identifies
    $P_B(W_j)$ with $P_B(W_0)$.
\end{proof}

\begin{theorem}[Common staircase state transferred between tensors]%
    \label{thm:peps_staircase_physical_operation_same_state}
    \lean{TNLean.PEPS.deformedRegionStateAssembled_staircasePhysicalOp_sameState}
    \leanok
    \uses{thm:peps_physical_operator_same_state_transfer,
        thm:peps_staircase_virtual_operation_physical_realization,
        thm:peps_staircase_virtual_operation_common_state,
        thm:peps_regionInteriorBondProd_staircaseWindow_eq}
    Let $L,K>0$, and choose staircase coordinates $(a,b)$ with
    $a+2L\leq\mathrm{width}$ and $2K\leq\mathrm{height}$. Let $A$ and $B$
    be translation invariant tensors with the same physical dimension and the
    same closed state. Suppose every cyclic $L\times K$ window of $A$ is
    injective and $D_A(g)>0$ for every edge $g$. For
    $X\in\MN{D_A(e)}$, let $O_j^A(X)$ be the physical operators obtained from
    the staircase insert family of $A$, and define inserts on the windows of
    $B$ by
    \begin{align}
        C_j^{A\to B}(X)(\mu)
         & =O_j^A(X)\bigl(T_{B,W_j}(\mu)\bigr).
        \notag
    \end{align}
    Then, for all $j,j'\in\mathbb{N}$,
    \begin{align}
        \Psi^{\mathrm{as}}_{B,W_j,C_j^{A\to B}(X)}
         & =\Psi^{\mathrm{as}}_{B,W_{j'},C_{j'}^{A\to B}(X)}.
        \notag
    \end{align}
    The source family has $0\leq j<L+K$. In the definition used here the
    truncated offsets stabilize at $W_{L+K-1}$ for $j\geq L+K-1$; the indices
    are not read modulo $L+K$. No equality between the edge-dependent bond
    dimensions of $A$ and $B$, and no positivity assumption on those of $B$,
    is required.
\end{theorem}
\begin{proof}
    \leanok
    For each $j$, (\ref{eq:peps_physical_op_same_state}) and the
    open-boundary realization of $O_j^A(X)$ give
    \begin{align}
        P_B(W_j)\Psi^{\mathrm{as}}_{A,W_j,C_j^A(X)}
         & =P_A(W_j)\Psi^{\mathrm{as}}_{B,W_j,C_j^{A\to B}(X)}.
        \notag
    \end{align}
    The staircase family of $A$ has one common deformed state. Translation
    invariance gives $P_A(W_j)=P_A(W_0)$ and
    $P_B(W_j)=P_B(W_0)$. Hence the right-hand side is independent of $j$.
    Positivity of the bond dimensions makes $P_A(W_0)$ nonzero, so it may be
    cancelled.
\end{proof}

\begin{definition}[The left staircase end operation $O_1^A(X)$]%
    \label{def:peps_staircase_O1_physical_operation}
    \lean{TNLean.PEPS.staircaseO1PhysicalOp}
    \leanok
    \uses{thm:peps_staircase_virtual_operation_physical_realization}
    For $X\in\MN{D_A(e)}$, define
    $O_1^A(X)=O_{L+K-1}^A(X)$ after identifying the last staircase
    window with the left end window.
\end{definition}

\begin{definition}[The underlying right staircase end operation $O_3^A(X)$]%
    \label{def:peps_staircase_O3_physical_operation}
    \lean{TNLean.PEPS.staircaseO3PhysicalOp}
    \leanok
    \uses{thm:peps_staircase_virtual_operation_physical_realization}
    For $X\in\MN{D_A(e)}$, define $O_3^A(X)=O_0^A(X)$ after identifying
    the first staircase window with the right end window.
    The relation (\ref{eq:peps_staircase_o3_cross_tensor_relation}) has the
    paper's transposed $O_3^{\mathsf T}$ orientation; no transpose is chosen
    on the abstract physical endomorphism space.
\end{definition}

\begin{theorem}[Cross-tensor virtual operation on the staircase edge]%
    \label{thm:peps_cross_tensor_virtual_operation_of_staircase}
    \lean{TNLean.PEPS.existsUnique_crossTensorVirtualOperation_of_staircasePhysicalOp_sameState}
    \leanok
    \uses{def:peps_staircase_O1_physical_operation,
        def:peps_staircase_O3_physical_operation}
    Let $L,K\geq2$, and choose the displayed non-wrapping horizontal
    staircase coordinates $(a,b)$ with
    \begin{align}
        1    & \leq a,
             & a+2L                 & \leq \mathrm{width},
             & b+2K-1               & \leq \mathrm{height},
        \notag                                              \\
        2L+1 & \leq \mathrm{width},
             & 2K+1                 & \leq \mathrm{height}.
        \notag
    \end{align}
    Suppose that $A$ and $B$ are translation invariant, generate the same
    closed state, every cyclic $L\times K$ window of either tensor is
    injective, and every bond dimension of either tensor is positive. For
    $X\in\MN{D_A(e)}$, there is a unique $Y\in\MN{D_B(e)}$ such that, for
    every external boundary configuration $\eta_{\mathrm L}$ and
    $\eta_{\mathrm R}$, and every reference-edge index $k$,
    \begin{align}
        O_1^A(X)\bigl(T_{B,W_{\mathrm L}}(k,\eta_{\mathrm L})\bigr)
         & =\sum_jY_{jk}T_{B,W_{\mathrm L}}(j,\eta_{\mathrm L}),
        \label{eq:peps_staircase_o1_cross_tensor_relation}       \\
        O_3^A(X)\bigl(T_{B,W_{\mathrm R}}(k,\eta_{\mathrm R})\bigr)
         & =\sum_jY_{kj}T_{B,W_{\mathrm R}}(j,\eta_{\mathrm R}).
        \label{eq:peps_staircase_o3_cross_tensor_relation}
    \end{align}
    Thus the relation (\ref{eq:peps_staircase_o3_cross_tensor_relation}) has the
    paper's $O_3^{\mathsf T}$ orientation.
    No equality between $D_A(e)$ and $D_B(e)$ is assumed or concluded, and
    this statement contains no multiplication or gauge assertion.
\end{theorem}
\begin{proof}
    \leanok
    \uses{thm:peps_staircase_physical_operation_same_state,
        thm:peps_staircase_pair_insert_eq_open,
        thm:peps_existsUnique_virtualOperation_of_horizontal_staircase_end_pair,
        thm:peps_regionInjectivityUnionClosure_of_overlap}
    The transferred physical inserts form one common $B$-state. The
    overlapping-union lemma supplies the union closure needed to compare
    consecutive staircase windows, and corner completion gives equality of
    the two end inserts with open boundary. At the two ends these inserts are
    precisely the actions of $O_3^A(X)$ and $O_1^A(X)$ on the genuine
    $B$-window tensors. Injectivity of those two windows and the two-end
    comparison give the displayed relations and the uniqueness of $Y$.
\end{proof}

\begin{theorem}[Composition of the canonical staircase end operations]%
    \label{thm:peps_staircase_end_physical_operation_multiplication}
    \lean{TNLean.PEPS.staircaseO1PhysicalOp_mul}
    \lean{TNLean.PEPS.staircaseO3PhysicalOp_mul}
    \leanok
    \uses{def:peps_staircase_O1_physical_operation,
        def:peps_staircase_O3_physical_operation}
    Let $L,K>0$, and choose staircase coordinates $(a,b)$ with
    $a+2L\leq\mathrm{width}$ and $2K\leq\mathrm{height}$. Suppose every
    cyclic $L\times K$ window of $A$ is injective. For matrices
    $X,Y\in\MN{D_A(e)}$, the canonical end operations satisfy
    \begin{align}
        O_1^A(XY) & =O_1^A(X)O_1^A(Y),
        \notag                         \\
        O_3^A(XY) & =O_3^A(Y)O_3^A(X).
        \notag
    \end{align}
    Thus the right relation preserves multiplication in ordinary order after
    passing to the orientation $(O_3^A)^{\mathsf T}$ used in Lemma~5.
\end{theorem}
\begin{proof}
    \leanok
    \uses{lem:peps_staircase_virtual_operation_end_inserts,
        thm:peps_lemma5_end_operation_multiplication}
    Split the virtual boundary of either end window into the reference-edge
    coordinate and the remaining external coordinates.  Let $E$ be this
    splitting map, and write $C_M$ for insertion of $M$ on the reference edge.
    On every fixed external fibre,
    \begin{align}
        C_M\bigl(E^{-1}(j,\eta)\bigr)
         & =\sum_i M_{ij}\,T_{A,W}\bigl(E^{-1}(i,\eta)\bigr).
        \notag
    \end{align}
    The chosen left inverse of the injective blocked tensor is a left inverse
    on the whole boundary-coordinate space, so this identity gives
    $O_W(M)O_W(N)=O_W(MN)$.  The first window carries $X^{\mathsf T}$; hence
    $(XY)^{\mathsf T}=Y^{\mathsf T}X^{\mathsf T}$ reverses the order of the
    underlying right operations.
\end{proof}

\begin{definition}[Canonical cross-tensor matrix assignment]%
    \label{def:peps_staircase_cross_tensor_virtual_operation}
    \lean{TNLean.PEPS.staircaseCrossTensorVirtualOperation}
    \lean{TNLean.PEPS.staircaseCrossTensorVirtualOperation_spec}
    \leanok
    \uses{thm:peps_cross_tensor_virtual_operation_of_staircase}
    Under the hypotheses of
    Theorem~\ref{thm:peps_cross_tensor_virtual_operation_of_staircase}, in the
    paper's notation the input and uniquely recovered matrices are
    called $X$ and $Y$.  Writing $Y_{A\to B}(X)$ only to distinguish the two
    tensor families, define
    \begin{align}
        Y_{A\to B}(X) & \in\MN{D_B(e)}
        \notag
    \end{align}
    to be the unique matrix satisfying
    (\ref{eq:peps_staircase_o1_cross_tensor_relation}) and
    (\ref{eq:peps_staircase_o3_cross_tensor_relation}) for every left and
    right external boundary configuration.  This definition makes no choice
    of an external boundary configuration.
\end{definition}

\begin{theorem}[Multiplicativity of the cross-tensor assignment]%
    \label{thm:peps_staircase_cross_tensor_virtual_operation_multiplication}
    \lean{TNLean.PEPS.staircaseCrossTensorVirtualOperation_mul}
    \leanok
    \uses{def:peps_staircase_cross_tensor_virtual_operation}
    Under the hypotheses of
    Theorem~\ref{thm:peps_cross_tensor_virtual_operation_of_staircase}, for all
    $X_1,X_2\in\MN{D_A(e)}$,
    \begin{align}
        Y_{A\to B}(X_1X_2)
         & =Y_{A\to B}(X_1)Y_{A\to B}(X_2).
        \notag
    \end{align}
\end{theorem}
\begin{proof}
    \leanok
    \uses{thm:peps_staircase_end_physical_operation_multiplication,
        thm:peps_lemma5_end_operation_multiplication}
    Set $Y_r=Y_{A\to B}(X_r)$ for $r=1,2$.  For every left external
    boundary configuration, the two $O_1$ relations compose in ordinary order:
    \begin{align}
        O_1^A(X_1)\Bigl(O_1^A(X_2)
        \bigl(T_{B,W_{\mathrm L}}(k,\eta_{\mathrm L})\bigr)\Bigr)
         & =\sum_i (Y_1Y_2)_{ik}\,
        T_{B,W_{\mathrm L}}(i,\eta_{\mathrm L}).
        \notag
    \end{align}
    For every right external boundary configuration, the underlying $O_3$
    operations compose in reverse order:
    \begin{align}
        O_3^A(X_2)\Bigl(O_3^A(X_1)
        \bigl(T_{B,W_{\mathrm R}}(k,\eta_{\mathrm R})\bigr)\Bigr)
         & =\sum_j (Y_1Y_2)_{kj}\,
        T_{B,W_{\mathrm R}}(j,\eta_{\mathrm R}).
        \notag
    \end{align}
    This is ordinary multiplication in the $O_3^{\mathsf T}$ orientation.  The
    preceding theorem identifies the left sides with the canonical operations
    for $X_1X_2$.  Both end relations therefore realize $Y_1Y_2$, and global
    uniqueness identifies this product with $Y_{A\to B}(X_1X_2)$.
\end{proof}

\begin{theorem}[End-window identity for one virtual operation]%
    \label{thm:peps_end_window_identity_staircase_virtual_operation}
    \lean{TNLean.PEPS.regionInsertedCoeff_endWindows_eq_staircaseVirtualOperation}
    \leanok
    \uses{lem:peps_assembleRegionSigma_restrict, lem:peps_reference_edge_left_endpoint_staircase_window, thm:peps_regionInsertedCoeff_eq_smul_edgeInsertedCoeff, thm:peps_regionInteriorBondProd_staircaseWindow_eq}
    Suppose the tensor is translation invariant and the staircase satisfies
    the stated size bounds.  Then for every $X\in\MN{D_B(e)}$ and every
    global physical configuration $\omega$,
    \begin{align}
        C_B(W_0,e,X^{\mathsf T};\omega|_{W_0},
        \omega|_{V\setminus W_0})
         & =
        C_B(W_{L+K-1},e,X;\omega|_{W_{L+K-1}},
        \omega|_{V\setminus W_{L+K-1}}).
        \notag
    \end{align}
\end{theorem}
\begin{proof}
    \leanok
    The region-to-edge identity sends both sides to the edge coefficient
    with $X$ on the same reference edge: the first boundary orientation
    transposes $X^{\mathsf T}$, while the last fixes $X$.  Translation
    invariance identifies the two non-boundary bond products.
\end{proof}


\begin{definition}[Window-region coefficient-identity witness]
    \label{def:peps_window_edge_coeff_identity_witness}
    \lean{TNLean.PEPS.windowEdgeCoeffIdentityWitness}
    \leanok
    \uses{def:peps_regionInsertedCoeff, thm:peps_staircase_end_reference_edge_boundary}
    Let $W$ be the left end window of a horizontal staircase with parameters
    $(L,K)$ and lower-left coordinate $(a,b)$, and let $e$ be its reference
    edge. Suppose that
    \begin{align}
        0<L,\qquad 0<K,\qquad
        1\le a,\qquad
        a+2L\le W_{\mathrm{torus}},\qquad
        b+2K-1\le H_{\mathrm{torus}}.
        \notag
    \end{align}
    Then $e\in\partial W$. Suppose also that the blocked tensors of $B$ on
    $W$ and $V\setminus W$ are injective, that all bond dimensions of $B$ are
    positive, and that the bond dimensions of $A$ and $B$ at $e$ are
    identified by $\iota_e:\MN{D_A(e)}\simeq\MN{D_B(e)}$. If
    $Z\in\GL(D_B(e))$ satisfies, for every matrix $M$ and physical
    configurations $\sigma,\tau$,
    \begin{align}
        C_A(W,e,M;\sigma,\tau)
        =
        C_B(W,e,Z\,\iota_e(M)\,Z^{-1};\sigma,\tau),
        \notag
    \end{align}
    then these objects form the coefficient-identity witness at $e$ with
    witness region $W$ and with both gauge entries equal to $Z$.
\end{definition}

\begin{definition}[Window-region witness from window injectivity]
    \label{def:peps_window_edge_coeff_identity_witness_of_hypotheses}
    \lean{TNLean.PEPS.windowEdgeCoeffIdentityWitness_of_hypotheses}
    \leanok
    \uses{def:peps_window_edge_coeff_identity_witness}
    Let $W$ be the left end window around the reference edge $e$, placed with
    lower-left coordinate $(a,b)$. Suppose that every cyclic $L\times K$
    window of $B$ is blocked-tensor injective, unions of injective regions are
    injective, all bond dimensions of $B$ are positive, and
    \begin{align}
        2\le L,\qquad 2\le K,\qquad
        2L+1\le W_{\mathrm{torus}},\qquad
        2K+1\le H_{\mathrm{torus}},
        \notag \\
        1\le a,\qquad
        a+2L\le W_{\mathrm{torus}},\qquad
        b+2K-1\le H_{\mathrm{torus}}.
        \notag
    \end{align}
    Suppose also that the bond dimensions of $A$ and $B$ at $e$ are
    identified by $\iota_e:\MN{D_A(e)}\simeq\MN{D_B(e)}$, and let
    $Z\in\GL(D_B(e))$. If the conjugation identity
    \begin{align}
        C_A(W,e,M;\sigma,\tau)
        =
        C_B(W,e,Z\,\iota_e(M)\,Z^{-1};\sigma,\tau)
        \notag
    \end{align}
    holds for every $M,\sigma,\tau$, then the coefficient-identity witness of
    Definition~\ref{def:peps_window_edge_coeff_identity_witness} is obtained
    for the region $W$. The hypotheses supply the injectivity of both $W$ and
    its complement at the minimal torus size.
\end{definition}

\paragraph{Window realization and witness read-off.}

\begin{lemma}[Window-level edge-insertion injectivity]%
    \label{lem:peps_windowRegionInsertedCoeff_injective}
    \lean{TNLean.PEPS.windowRegionInsertedCoeff_injective}
    \leanok
    \uses{thm:peps_regionInsertedCoeff_injective,
        lem:peps_horizontal_staircase_window_injectivity,
        thm:peps_horizontal_staircase_reference_edge_crossing}
    Let $B$ be a PEPS tensor on a torus whose periods $W,H$ satisfy
    $1<W$ and $1<H$.  Assume the one-orientation $L\times K$
    window-injectivity hypotheses for $B$, union closure for injective
    regions, positive bond dimensions, and
    \[
        2\le L,\qquad 2\le K,\qquad 1\le a,\qquad
        a+2L\le W,\qquad b+2K-1\le H,
    \]
    together with
    \[
        2L+1\le W,\qquad 2K+1\le H.
    \]
    Let $W_{\mathrm L}$ be the left end window and let $e$ be the reference
    edge.  If $M,M'\in\MN{D_B(e)}$ satisfy
    \[
        C_B(W_{\mathrm L},e,M;\sigma,\tau)
        =
        C_B(W_{\mathrm L},e,M';\sigma,\tau)
    \]
    for every physical configuration $\sigma$ on $W_{\mathrm L}$ and $\tau$
    on $V\setminus W_{\mathrm L}$, then $M=M'$.
\end{lemma}
\begin{proof}
    \leanok
    The window-injectivity hypotheses give injectivity of the blocked tensor
    map on $W_{\mathrm L}$.  The single-window complement is also injective at
    the displayed size.  Since $e\in\partial W_{\mathrm L}$,
    Theorem~\ref{thm:peps_regionInsertedCoeff_injective} applied to the
    region $W_{\mathrm L}$ gives $M=M'$.
\end{proof}

\begin{theorem}[Window witness from a region insertion transfer]%
    \label{thm:peps_windowEdgeCoeffIdentityWitness_of_transfer}
    \lean{TNLean.PEPS.exists_windowEdgeCoeffIdentityWitness_of_transfer}
    \leanok
    \uses{def:peps_region_insertion_transfer,
        def:peps_window_edge_coeff_identity_witness,
        def:peps_window_edge_coeff_identity_witness_of_hypotheses}
    Let $A$ and $B$ be PEPS tensors on a torus whose periods $W,H$ satisfy
    $1<W$ and $1<H$.  Assume the one-orientation $L\times K$
    window-injectivity hypotheses for both tensors, union closure for
    injective regions for both tensors, positive bond dimensions, and
    \[
        2\le L,\qquad 2\le K,\qquad 1\le a,\qquad
        a+2L\le W,\qquad b+2K-1\le H,
    \]
    together with
    \[
        2L+1\le W,\qquad 2K+1\le H.
    \]
    Let $W_{\mathrm L}$ be the left end window and let $e$ be the reference
    edge.  If $T_e:\MN{D_A(e)}\to\MN{D_B(e)}$ is a region insertion transfer
    on $W_{\mathrm L}$ at $e$, then there exist
    $Z,Z_{\mathrm{ref}}\in\GL(D_B(e))$ and a bond-dimension identification
    $\iota_e:\MN{D_A(e)}\simeq\MN{D_B(e)}$ such that the coefficient-identity
    witness at $e$ is nonempty.  The coefficient identity has the form
    \[
        C_A(W_{\mathrm L},e,M;\sigma,\tau)
        =
        C_B(W_{\mathrm L},e,Z\,\iota_e(M)\,Z^{-1};\sigma,\tau).
    \]
\end{theorem}
\begin{proof}
    \leanok
    The transfer determines a gauge $Z$ by the region insertion algebra:
    \[
        T_e(M)=Z\,\iota_e(M)\,Z^{-1}.
    \]
    Its coefficient-matching property gives the displayed identity.  The
    window and complement injectivity hypotheses are those required by
    Definition~\ref{def:peps_window_edge_coeff_identity_witness_of_hypotheses},
    which yields the witness with $Z_{\mathrm{ref}}=Z$.
\end{proof}

\begin{theorem}[Window witness from single-vertex realization]%
    \label{thm:peps_windowEdgeCoeffIdentityWitness_of_realizes}
    \lean{TNLean.PEPS.exists_windowEdgeCoeffIdentityWitness_of_realizes}
    \leanok
    \uses{def:peps_region_transfer_realizes,
        def:peps_region_insertion_transfer_of_realizes,
        thm:peps_windowEdgeCoeffIdentityWitness_of_transfer}
    Let $A$ and $B$ be PEPS tensors on a torus whose periods $W,H$ satisfy
    $1<W$ and $1<H$.  Assume the one-orientation $L\times K$
    window-injectivity hypotheses for both tensors, union closure for
    injective regions for both tensors, matched bond dimensions, equality of
    the PEPS state, positive bond dimensions, and
    \[
        2\le L,\qquad 2\le K,\qquad 1\le a,\qquad
        a+2L\le W,\qquad b+2K-1\le H,
    \]
    together with
    \[
        2L+1\le W,\qquad 2K+1\le H.
    \]
    Let $W_{\mathrm L}$ be the left end window and let $e$ be the reference
    edge.  Suppose the endpoint tensor of each family at the in-region
    endpoint of $e$ has linearly independent columns, and suppose the boundary
    matrix transfer is realized on $W_{\mathrm L}$ in both directions.  Then
    there exist $Z,Z_{\mathrm{ref}}\in\GL(D_B(e))$ and a bond-dimension
    identification $\iota_e:\MN{D_A(e)}\simeq\MN{D_B(e)}$ such that the
    coefficient-identity witness at $e$ is nonempty.
\end{theorem}
\begin{proof}
    \leanok
    The two realized boundary transfers determine a region insertion transfer:
    \[
        C_A(W_{\mathrm L},e,M;\sigma,\tau)
        =
        C_B(W_{\mathrm L},e,T_e(M);\sigma,\tau),
        \qquad
        T_e(MM')=T_e(M)T_e(M').
    \]
    Theorem~\ref{thm:peps_windowEdgeCoeffIdentityWitness_of_transfer} then
    gives the two gauges and the nonempty coefficient-identity witness.
\end{proof}

\paragraph{Window witness from bond locality.}

\begin{theorem}[Forward window coefficient transfer from bond locality]
    \label{thm:peps_windowCoeffTransferAB_of_bondLocal}
    \lean{TNLean.PEPS.windowCoeffTransferAB_of_bondLocal}
    \leanok
    \uses{def:peps_regionInsertedCoeff,
        lem:peps_horizontal_staircase_window_injectivity,
        thm:peps_horizontal_staircase_reference_edge_crossing,
        thm:peps_region_transfer_coeff_double_factorization,
        thm:peps_region_transfer_coeff_incident_form}
    Let $A$ and $B$ be PEPS tensors on a torus whose periods
    $W_{\mathrm{torus}},H_{\mathrm{torus}}$ satisfy $1<W_{\mathrm{torus}}$
    and $1<H_{\mathrm{torus}}$.  Assume the one-orientation $L\times K$
    window-injectivity hypotheses and union closure for both tensors, equality
    of the PEPS state, matched bond dimensions, positive bond dimensions, and
    \[
        2\le L,\qquad 2\le K,\qquad 1\le a,\qquad
        a+2L\le W_{\mathrm{torus}},\qquad
        b+2K-1\le H_{\mathrm{torus}},
    \]
    together with
    \[
        2L+1\le W_{\mathrm{torus}},\qquad
        2K+1\le H_{\mathrm{torus}}.
    \]
    Let $W_{\mathrm L}$ be the left end window and let $e$ be the reference
    edge.  Suppose the $A$-to-$B$ transfer kernel on $W_{\mathrm L}$ at $e$ is
    incident along the boundary leg $e$.  Then for every
    $M\in\MN{D_A(e)}$ there is a matrix $N\in\MN{D_B(e)}$ such that
    \[
        C_A(W_{\mathrm L},e,M;\sigma,\tau)
        =
        C_B(W_{\mathrm L},e,N;\sigma,\tau)
    \]
    for every physical configuration $\sigma$ on $W_{\mathrm L}$ and $\tau$
    on $V\setminus W_{\mathrm L}$.
\end{theorem}
\begin{proof}
    \leanok
    The window hypotheses give injectivity of the blocked tensor maps on
    $W_{\mathrm L}$ for $A$ and $B$, while union closure gives injectivity on
    the complements.  The coefficient-transfer theorem for the blocked region
    and its complement rewrites the first tensor's inserted coefficient as a
    double contraction through the second tensor.  The incident form of the
    transfer kernel along $e$ turns this double contraction into the single
    region insertion $C_B(W_{\mathrm L},e,N;\sigma,\tau)$.
\end{proof}

\begin{theorem}[Backward window coefficient transfer from bond locality]
    \label{thm:peps_windowCoeffTransferBA_of_bondLocal}
    \lean{TNLean.PEPS.windowCoeffTransferBA_of_bondLocal}
    \leanok
    \uses{def:peps_regionInsertedCoeff,
        lem:peps_horizontal_staircase_window_injectivity,
        thm:peps_horizontal_staircase_reference_edge_crossing,
        thm:peps_region_transfer_coeff_double_factorization,
        thm:peps_region_transfer_coeff_incident_form}
    Under the hypotheses of
    Theorem~\ref{thm:peps_windowCoeffTransferAB_of_bondLocal}, suppose instead
    that the $B$-to-$A$ transfer kernel on $W_{\mathrm L}$ at $e$ is incident
    along the boundary leg $e$.  Then for every $N\in\MN{D_B(e)}$ there is a
    matrix $M\in\MN{D_A(e)}$ such that
    \[
        C_B(W_{\mathrm L},e,N;\sigma,\tau)
        =
        C_A(W_{\mathrm L},e,M;\sigma,\tau)
    \]
    for every physical configuration $\sigma$ on $W_{\mathrm L}$ and $\tau$
    on $V\setminus W_{\mathrm L}$.
\end{theorem}
\begin{proof}
    \leanok
    Apply the forward theorem with the two tensors interchanged.  The same
    window and complement injectivity hypotheses are available in the reversed
    direction, and the incident transfer kernel now supplies the matrix $M$ on
    the $A$ boundary bond.
\end{proof}

\begin{theorem}[Window witness from bond locality and multiplicativity]
    \label{thm:peps_windowEdgeCoeffIdentityWitness_of_bondLocal}
    \lean{TNLean.PEPS.exists_windowEdgeCoeffIdentityWitness_of_bondLocal}
    \leanok
    \uses{thm:peps_windowCoeffTransferAB_of_bondLocal,
        thm:peps_windowCoeffTransferBA_of_bondLocal,
        def:peps_coeff_transfer_map,
        thm:peps_region_edge_gauge_of_coeff_transfer,
        def:peps_window_edge_coeff_identity_witness,
        def:peps_window_edge_coeff_identity_witness_of_hypotheses}
    Assume the hypotheses of
    Theorem~\ref{thm:peps_windowCoeffTransferAB_of_bondLocal} and its
    backward analogue
    Theorem~\ref{thm:peps_windowCoeffTransferBA_of_bondLocal}.  Let
    $T_e:\MN{D_A(e)}\to\MN{D_B(e)}$ be the coefficient transfer chosen from
    the forward theorem.  If
    \[
        T_e(MM')=T_e(M)T_e(M')
    \]
    for all $M,M'\in\MN{D_A(e)}$, then there exist
    $Z,Z_{\mathrm{ref}}\in\GL(D_B(e))$ and a bond-dimension identification
    $\iota_e:\MN{D_A(e)}\simeq\MN{D_B(e)}$ such that the window-region
    coefficient-identity witness at $e$ is nonempty.  In particular the
    coefficient identity has the form
    \[
        C_A(W_{\mathrm L},e,M;\sigma,\tau)
        =
        C_B(W_{\mathrm L},e,Z\,\iota_e(M)\,Z^{-1};\sigma,\tau).
    \]
\end{theorem}
\begin{proof}
    \leanok
    The two bond-locality hypotheses give the forward and backward coefficient
    transfers by
    Theorems~\ref{thm:peps_windowCoeffTransferAB_of_bondLocal}
    and~\ref{thm:peps_windowCoeffTransferBA_of_bondLocal}.  With the displayed
    multiplicativity, the coefficient-transfer construction gives a conjugating
    gauge $Z$ and the coefficient identity on $W_{\mathrm L}$.  The window and
    complement injectivity hypotheses then supply the window-region
    coefficient-identity witness with $Z_{\mathrm{ref}}=Z$.
\end{proof}

